\begin{theorem} \label{thm:leadingonesexpo}
The expected runtime of the \new{$(1+1)$~Opt-IA$_{\textsc{FirstBest}}$}  %with {\expoHD } 
with {\linHD } \new{and any $c\leq 1$} for optimising
\textsc{LeadingOnes} is $\Theta(n^3)$ with $\tau=\Omega(n^{1+\epsilon})$ for any arbitrary 
constant $\epsilon>0$.
\end{theorem}
\begin{proof}
	
Fix a level~$\ell$ with $n/3\le \ell\le 2n/3$. For the current individual $x$ with $\mathrm{LO}(x)=\ell$, we have $x_1=\cdots=x_{\ell}=1$ and $x_{\ell+1}=0$, and the remaining $m:=n-\ell-1$ positions form the tail. In each iteration, the first flipped bit is uniformly distributed over $\{1,\dots,n\}$; therefore the probability that the first flip hits the improving bit $\ell+1$ is $1/n$, the probability that it hits the prefix is $\ell/n$, and the probability that it hits the tail is $m/n$. Let $T$ be the iteration on which the improving bit is first hit. Then $T$ is geometrically distributed with parameter~$1/n$, so $\mathbb{E}[T]=n$. Before this improving iteration occurs, the expected number of tail-first iterations is $(m/n)(\mathbb{E}[T]-1)=m=\Theta(n)$. Each such iteration flips a uniformly random tail bit. By standard coupon–collector arguments, after $\Theta(n)$ such tail flips a constant fraction of the $m$ tail positions has been flipped at least once with constant probability. Thus, with constant probability, the Hamming distance between $x$ and the best is $\Theta(m)=\Theta(n)$, and therefore the IPH mutation potential satisfies $M=\Theta(n)$ at this level.
	
	Consider now a subsequent iteration at the same level while $M\ge \alpha n$ for some constant $\alpha>0$. If the first flip hits the prefix (an event of probability at least $\ell/n\ge 1/3$), then, because mutation proceeds without replacement, the lost prefix bit cannot be repaired within the same iteration. Consequently, no acceptance can occur in that iteration, and the algorithm executes the entire mutation budget $M$, rejecting at the end. Hence the expected number of fitness evaluations in such an iteration is at least $\alpha n$. Since this happens with probability at least $1/3$, the unconditional expected cost of one iteration while $M=\Theta(n)$ is $\Omega(n)$.
	
	The probability that an iteration yields an improvement is exactly $1/n$, because only a first flip that hits position $\ell+1$ increases the \textsc{LeadingOnes} value. Therefore the expected number of iterations required to advance from level $\ell$ to $\ell+1$ is $1/(1/n)=n$. Combining this with the previous observation, the expected number of fitness evaluations spent at level $\ell$ satisfies $\Omega(n)\cdot n=\Omega(n^2)$.
	There are $\Theta(n)$ levels $\ell$ in the interval $[n/3,2n/3]$, and each such level contributes $\Omega(n^2)$ expected evaluations. Thus the total expected optimisation time is $\Omega(n^3)$. Finally, stochastic ageing can only increase the number of evaluations and therefore cannot invalidate the lower bound. Since the function has $n$ levels, $1/n$ probability of improvement and $M\leq n$, the expected runtime is bounded from above by $O(n^3)$ as well.
	
	
\end{proof}
